\section{Auswertung der Beanspruchung}
\label{sec:auswertung_beanspruchung}

Nach jedem Lauf wird die Ergebnisdatenbank der Stellung unter dem Python von Abaqus ausgelesen. Eine Felddatei hält je Laststufe die Reaktionen der Referenzpunkte, die geschlossenen Kontaktknoten der Radflanke mit ihrer Pressung, die knotengemittelten Spannungen jeder Fußrundungsfläche und die größte Kopfverschiebung fest. Von den Fußrundungen stehen darin die Knoten der Breitenmitte vollständig, je Schicht über der Zahnbreite die größte Zug- und die betragsgrößte Druckhauptspannung. Eine zweite Datei enthält den knotengemittelten Spannungstensor aller Knoten des Rades in dessen Teilesystem, in das die Instanzdrehung und ein lokales Orientierungssystem der Stufen W3 und W4 zurückgerechnet werden. Derselbe Knoten trägt so in jeder Stellung dasselbe Bezugssystem, und seine Lastgeschichte lässt sich Komponente für Komponente zusammensetzen. Gemittelt wird über die Elementknotenwerte der beteiligten Elemente, Invarianten entstehen erst aus dem gemittelten Tensor.

Für die Kurven über dem Wälzweg wird jede Stellung auf ihre Größtwerte verdichtet. Festgehalten werden je Zahnhälfte der größte Kontaktdruck und die Normalkraft, je Fußrundung die größte Zughauptspannung $\sigma_{1}$ und die betragsgrößte Druckhauptspannung $\sigma_{3}$ mit ihrer Schicht, dazu Kopfverschiebung und Ritzeldrehung. Aufgereiht nach dem Abstand vom Eingriffsbeginn A ergeben die Stellungen eines Ordners den Verlauf dieser Größen als Größtwerte über die Zahnbreite. Die Normpunkte A bis E sind im Stellungsgitter stets enthalten und werden als senkrechte Linien eingetragen; je Größe werden das Maximum, sein Ort ab A und die Werte an den fünf Punkten tabelliert. Ein abgebrochener Schritt liefert nur sein letztes konvergiertes Inkrement zu einer kleineren Last und bleibt außerhalb der Kurven.

Bewertet werden die Spannungen mit dem vereinfachten Festigkeitsnachweis für thermoplastische Kunststoffbauteile nach Stommel et al. \cite[Kap.~6]{Sto18}, einem Auslastungsnachweis je FE-Knoten. Weil der Ort der höchsten Beanspruchung nicht der Ort des höchsten Auslastungsgrads sein muss \cite[S.~313]{Sto18}, wird jeder Knoten des Rades bewertet. Seine Lastgeschichte besteht aus dem unbelasteten Zustand und den 99 Stellungen einer Serie und Laststufe; ein Fußknoten erfährt darin den Zug unter dem eigenen Zahn und den Druck, während der Nachbarzahn trägt, und gilt ohne weitere Annahme als wechselnd beansprucht. Das Buch verlangt für die Spannungsermittlung ein isotropes Werkstoffmodell, auch bei kurzfaserverstärkten Kunststoffen, und schließt anisotrope Modelle ausdrücklich aus \cite[S.~343 bis 344]{Sto18}; im Ermüdungsnachweis werden die Spannungen stets linear-elastisch ermittelt \cite[S.~362]{Sto18}. Grundlage des Nachweises wie veröffentlicht ist daher die Stufe W1, daneben wird W2 ausgewertet, weil der statische Nachweis bei kurzzeitiger Beanspruchung auch elasto-plastisch ermittelte Spannungen zulässt \cite[S.~343]{Sto18}.

Im statischen Nachweis werden negative Hauptspannungen zuerst durch den Druckfestigkeitsfaktor $f_{\sigma} = 1{,}3$ geteilt \cite[S.~357]{Sto18}. In der Vergleichsspannung vermittelt dann die Steuergröße $q$, die aus dem Schubfestigkeitsfaktor $f_{\tau}$ der Werkstoffklasse folgt, zwischen von-Mises-Spannung und Normalspannungshypothese \cite[S.~322 und 357]{Sto18}. Der Auslastungsgrad setzt die zu $\sigma_{\mathrm{V}}$ auf der Betriebsisotherme abgelesene Vergleichsdehnung ins Verhältnis zur zulässigen Dehnung \cite[S.~355]{Sto18}:
\begin{equation}
\sigma_{\mathrm{V}} = q\,\sigma_{\mathrm{V,NH}} + (1 - q)\,\sigma_{\mathrm{V,Mises}} , \qquad
a_{\mathrm{stat}} = \frac{\varepsilon_{\mathrm{V}}(\sigma_{\mathrm{V}})}{\varepsilon_{\mathrm{zul}}(h)/S} .
\label{eq:a_stat}
\end{equation}
Die zulässige Dehnung hängt vom Mehraxialitätsgrad $h$ nach Hencky ab, dem Verhältnis von hydrostatischer zu von-Mises-Spannung; bis $h = 1/3$ gilt die obere Dehnungsgrenze, darüber fällt $\varepsilon_{\mathrm{zul}}$ nichtlinear auf die untere ab \cite[S.~325 bis 326]{Sto18}. Beide Grenzen liefert das Nomogramm in \figref{fig:nomogramm_qs}. Der obere und der untere Abminderungsfaktor der Klasse und Lastart skalieren dort den statischen Werkstoffkennwert $K$ der Raumtemperaturkurve, der Schnitt mit der Referenzkurve legt die Dehnungen fest \cite[S.~344 bis 348]{Sto18}. Diese gelten als temperaturunabhängig und werden senkrecht auf die Isotherme der Betriebstemperatur übertragen, von der die statische Bauteilfestigkeit $\sigma_{\mathrm{stat}}$ abgelesen wird \cite[S.~349 bis 350]{Sto18}. Dem Wert $A = 1{,}0$ entspricht die Streckgrenze der Referenzkurve \cite[S.~346 und 349]{Sto18}; weil das konditionierte PA46-GF30 keinen ausgeprägten Streckpunkt zeigt, dient das Maximum der Kurve als $K$. Der Zustand mit dem größten Auslastungsgrad ist der maßgebende Zustand des Knotens und gibt dem zyklischen Nachweis $h$, $\sigma_{\mathrm{stat}}$ und die Hauptrichtungen vor; liegt $\sigma_{\mathrm{V}}$ jenseits des Endes der Isotherme, wird der Auslastungsgrad als unendlich ausgewiesen.

Der zyklische Nachweis vergleicht die Spannungsamplitude mit der Bauteil-Betriebsfestigkeit
\begin{equation}
\sigma_{\mathrm{B}} = K_{\mathrm{B}}\,K_{\mathrm{M}}\,\sigma_{\mathrm{W}} \qquad \text{mit} \qquad \sigma_{\mathrm{W}} = A_{\mathrm{W}}\,K ,
\label{eq:betriebsfestigkeit}
\end{equation}
deren Wechselfestigkeit aus dem zyklischen Abminderungsfaktor für die Grenzlastspielzahl $N_{\mathrm{G}} = 10^{7}$ hervorgeht, für kurzfaserverstärkte Klassen ein fester Wert \cite[S.~360 bis 363]{Sto18}. Den Mittelspannungseinfluss trägt $K_{\mathrm{M}}$ aus dem Haigh-Diagramm in \figref{fig:haigh_qs_80C}. Dort läuft die Goodman-Gerade von $\sigma_{\mathrm{W}}$ zur Streckgrenze $\sigma_{\mathrm{Y}}$ der Betriebsisotherme mit der Mittelspannungsempfindlichkeit $M = \sigma_{\mathrm{W}}/\sigma_{\mathrm{Y}}$ und wird durch $\sigma_{\mathrm{stat}}$ begrenzt \cite[S.~365]{Sto18}. Auf welchem Weg die Beanspruchung die Grenzlinie erreicht, legt der Überlastungsfall fest \cite[S.~366 bis 368]{Sto18}; gewählt ist Fall F2 mit konstantem Spannungsverhältnis $R$. Für Lastspielzahlen $N$ unterhalb von $N_{\mathrm{G}}$ folgt der Betriebsfestigkeitsfaktor einer synthetischen Wöhlerlinie von $\sigma_{\mathrm{stat}}$ bei $N = 1$ zu $\sigma_{\mathrm{W}}$ bei $N_{\mathrm{G}}$, gebildet mit den auf die Betriebstemperatur übertragenen Faktoren \cite[S.~369 bis 371]{Sto18}:
\begin{equation}
K_{\mathrm{M}} = \frac{1}{1 + M\,\sigma_{\mathrm{m}}/\sigma_{\mathrm{a}}} , \qquad
K_{\mathrm{B}} = \left(\frac{N_{\mathrm{G}}}{N}\right)^{1/k} , \qquad
k = \frac{\lg N_{\mathrm{G}}}{\lg\left(A_{\mathrm{S}}/A_{\mathrm{W}}\right)} .
\label{eq:k_m_k_b}
\end{equation}
Der Fall F2 ist gewählt, weil das Zahnrad mit dem Drehmoment einen einzigen Lastparameter besitzt und jede Spannung des linearen Modells mit ihm skaliert. Der Überlastungsfall beschreibt damit den Lastpfad und bleibt als einzige Festlegung für alle Knoten gleich, während Amplitude, Mittelspannung, $R$, $h$ und $K_{\mathrm{M}}$ aus der Geschichte jedes Knotens entstehen. Gedeckelt bleibt die Betriebsfestigkeit durch $\sigma_{\mathrm{stat}} - \sigma_{\mathrm{m}}$ \cite[S.~372]{Sto18}, angesetzt mit der statischen Grenze des maßgebenden Zustands und dem Betrag der Mittelspannung auf beiden Seiten. Die Lebensdauer eines Knotens ist die Lastspielzahl, bei der der Vergleichsauslastungsgrad auf dieser Linie eins erreicht; unterhalb von $10^{4}$ Lastspielen gilt das Verfahren nicht \cite[S.~372]{Sto18}, oberhalb von $N_{\mathrm{G}}$ ist mit weiterem Abfall zu rechnen \cite[S.~328]{Sto18}. \Note{PH: Empfindlichkeit der Überlastungsfälle F1 bis F4 an den kritischen Knoten (E2), noch offen.}

Geführt wird der zyklische Nachweis komponentenweise, weil sich im mehraxialen Fall keine Vergleichsspannung bilden lässt \cite[S.~368 und 373]{Sto18}; die Komponenten sind die Hauptnormalspannungen unter der Annahme nicht drehender Hauptrichtungen \cite[S.~361]{Sto18}. Dafür werden die Hauptrichtungen des maßgebenden Zustands festgehalten und alle Zustände der Geschichte auf sie projiziert, so dass eine Komponente ihr Vorzeichen behält, wenn sie zwischen Zug und Druck wechselt. Je Richtung $i$ folgen Amplitude $\sigma_{\mathrm{a},i}$, Mittelspannung und $K_{\mathrm{M}}$, der Einzelauslastungsgrad $a_{\sigma,i} = \sigma_{\mathrm{a},i}/\sigma_{\mathrm{B}}$ und der Vergleichsauslastungsgrad \cite[S.~374]{Sto18}
\begin{equation}
a_{\mathrm{B,V}} = q\,a_{\mathrm{B,NH}} + (1 - q)\,a_{\mathrm{B,Mises}} \qquad \text{mit} \qquad a_{\mathrm{B,NH}} = \max\left(a_{\sigma,1}, a_{\sigma,2}, a_{\sigma,3}\right)
\label{eq:a_bv}
\end{equation}
mit $a_{\mathrm{B,Mises}}$ als von-Mises-Verknüpfung der drei Einzelauslastungsgrade. Als Vergleichswege sind die vorzeichenbehaftete von-Mises-Spannung und die komponentenweise Bildung von Amplitude und Mittelwert hinterlegt \cite[S.~332]{Sto18}. \Note{PH: Seiten der Gleichungen (6.10) bis (6.12) prüfen}

Zahlenwerte setzt die Werkstoffklasse. Für die Abminderungsfaktoren zählt das konditionierte PA46-GF30 zu den teilkristallinen verstärkten Thermoplasten mit $A_{\mathrm{S}}$ von $0{,}55$ bis $0{,}70$ für die einmalige Lastart und $A_{\mathrm{W}} = 0{,}25$ \cite[S.~346 und 362]{Sto18}. Für die Versagenshypothese fallen feuchte hydrophile verstärkte Thermoplaste nach der Fußnote der Klassentabelle in die duktile Spalte mit $q = 0$ \cite[S.~323]{Sto18}. Matscape schätzt den Versagenstyp der Karten aus den Versuchsdaten ebenfalls als duktil mit $f_{\tau} = 0{,}577$; dieser Wert wird übernommen, womit die Vergleichsspannung zur druckkorrigierten von-Mises-Spannung wird. Im zyklischen Nachweis entfällt bei der Steuergröße die Druckkorrektur \cite[S.~374]{Sto18}.

Die Stufen W3 und W4 rechnen mit der anisotropen CONVERSE-Karte, die das veröffentlichte Verfahren ausschließt. Als eigene Erweiterung dieser Arbeit ersetzt die Hill-Vergleichsspannung im Faserkoordinatensystem die von-Mises-Spannung in \eqref{eq:a_stat},
\begin{equation}
\sigma_{\mathrm{V,Hill}}^{2} = F\left(\sigma_{22} - \sigma_{33}\right)^{2} + G\left(\sigma_{33} - \sigma_{11}\right)^{2} + H\left(\sigma_{11} - \sigma_{22}\right)^{2} + 2L\,\sigma_{23}^{2} + 2M\,\sigma_{13}^{2} + 2N\,\sigma_{12}^{2} ,
\label{eq:hill}
\end{equation}
deren sechs Konstanten aus den Fließspannungsverhältnissen $R_{ij}$ der Karte folgen. Das ist das quadratische Hill-Potential, das Abaqus für die von CONVERSE geschriebene Karte auswertet \cite{AbaqusHill25}. Die Festigkeitsseite der $0$°-Kurve bleibt unverändert, weil die Verhältnisse der Karte jeden Orientierungszustand auf eben diesen Stab beziehen; die Orientierung wirkt allein über die Vergleichsspannung. Je Knoten wird das Faserkoordinatensystem vorzeichenfrei aus den Elementorientierungen gemittelt, der Orientierungszustand aus den Feldvariablen gelesen und die Hill-Zeile über das Zustandsgitter der Karte bei Betriebstemperatur linear interpoliert. Die druckkorrigierten Hauptspannungen werden auf ihre Richtungen zurückgelegt und in das Faserkoordinatensystem gedreht, bevor \eqref{eq:hill} ausgewertet wird; Mehraxialitätsgrad und Vorzeichen der Geschichte bleiben als Invarianten bezugsfrei, und der zyklische Nachweis läuft mit der vorzeichenbehafteten Hill-Spannung. Offen ist die Klasse der Abminderungsfaktoren. Das Buch erfasst die Festigkeitsminderung quer zur Faser pauschal über den niedrigen Faktor der verstärkten Klasse \cite[S.~345 bis 346]{Sto18}; mit aufgelöster Anisotropie zählte dieselbe Minderung zweimal. Voreingestellt sind unter Hill deshalb die Faktoren der teilkristallinen unverstärkten Klasse, wobei der niedrige Faktor der verstärkten Klasse nach dem Buch auch die geringe plastische Reserve abdeckt \cite[S.~345]{Sto18}. \Note{PH: Klassenwahl unter Hill (verstärkt gegen duktil unverstärkt), Entscheidung des Nutzers.}

Geprüft wurde die Umsetzung auf drei Wegen. Die Gleichungen des sechsten Kapitels liegen als einzeln nachvollziehbare Funktionen vor; ihre Prüffälle sind die Zahlenbeispiele des Buches zu Nomogramm und Temperaturübertragung \cite[S.~346 bis 350]{Sto18}. Die Werkstoffseite wurde gegen Matscape geprüft, das je Karte und Temperatur die statische Grenze, die Wechselfestigkeit und die Mittelspannungsempfindlichkeit mit denselben Abminderungsfaktoren ausgibt; dieselben Größen werden aus den Rohkurven nachgerechnet. Den Nachweis selbst bestätigt S-Life Plastics, die Software, auf die das Buch für die praktische Anwendung verweist \cite[S.~375]{Sto18}. Dafür wurde die Stellung C des eingefrorenen Netzes mit W1 in S-Life Plastics 4.0.0 auf der gesamten Radinstanz bewertet. Eingestellt waren die drei Laststufen als Lastfälle, $R = 0$, Überlastungsfall F2, Grenzlastspielzahl $10^{7}$, die quasistatische Karte bei $80$\,°C und ungemittelte Spannungen. Fünf benannte Knoten wurden verglichen, zwei in der Zugfußrundung, einer in der Druckfußrundung, einer am Zahnkopf und einer unter der Oberfläche; der letzte erhält keinen Bericht, weil S-Life nur Oberflächenknoten bewertet. Die Berichte je Knoten und Laststufe wurden über ihre Textebene in eine Datendatei übertragen und neben den eigenen Nachweis mit der Lastgeschichte aus unbelastetem Zustand und Stellung C gestellt. Verglichen werden die Hauptspannungen, weil beide Programme den Tensor in verschiedenen Systemen führen, danach die Glieder beider Nachweise; Unterschiede in den Konventionen werden im Ergebnisteil erklärt. \Note{PH: Verweis auf den Abschnitt des Nachweises in Kapitel 5} Die Einstellungen des Nachweises stehen in \tabref{tab:nachweis_einstellungen}.

\begin{table}[htbp]
    \caption{Einstellungen des Nachweises nach Stommel et al. für das Rad aus PA46-GF30}
    \label{tab:nachweis_einstellungen}
    \centering
    \begin{tabular}{L{0.28\textwidth} L{0.64\textwidth}}
        \hline
        Einstellung & Wert \\
        \hline
        Bewertete Knoten & alle Knoten des Rades; berichtet werden die beiden Fußrundungen des mittleren Zahns und das globale Maximum \\
        Werkstoffstufen & W1 nach dem Verfahren wie veröffentlicht, W2 daneben; W3 und W4 mit der Hill-Erweiterung \\
        Lastgeschichte & unbelasteter Zustand und 99 Stellungen je Serie und Laststufe (8, 12 und $16$\,Nm) \\
        Spannungen & Knotenmittel der Elementknotenwerte des Tensors im Teilesystem des Rades \\
        Referenzkurve & Zugprobe $0$° (Fließrichtung), quasistatische Karte in den Rohwerten von Matscape, technische Spannung, $23$\,°C; Betriebsisotherme $80$\,°C \\
        Werkstoffkennwert $K$ & Maximum der Kurve bei $23$\,°C (kein ausgeprägter Streckpunkt) \\
        Lastart statisch & einmalig, $A_{\mathrm{S}}$ $0{,}55$ bis $0{,}70$ (teilkristallin, verstärkt) \\
        Zyklischer Abminderungsfaktor & $A_{\mathrm{W}} = 0{,}25$ (teilkristallin, verstärkt), $N_{\mathrm{G}} = 10^{7}$ \\
        Versagenshypothese & duktil, $q = 0$, $f_{\tau} = 0{,}577$ (Matscape) \\
        Druckfestigkeitsfaktor & $f_{\sigma} = 1{,}3$ für negative Hauptspannungen \\
        Mehraxialität & $h$ nach Hencky je Zustand; $h = 0$ im unbelasteten Zustand \\
        Zyklischer Weg & Hauptrichtungen des maßgebenden Zustands, über die Geschichte festgehalten \\
        Überlastungsfall & F2 ($R$ konstant), für alle Knoten gleich \\
        Geforderte Lastspielzahl & $N = N_{\mathrm{G}}$ ($K_{\mathrm{B}} = 1$); Lebensdauer über die synthetische Wöhlerlinie nach \eqref{eq:k_m_k_b} \\
        Statische Kappe & $\sigma_{\mathrm{stat}}(h) - |\sigma_{\mathrm{m}}|$ des maßgebenden Zustands, beidseitig \\
        Sicherheitsfaktor & $S = 1$ \\
        Klasse unter Hill & teilkristallin, unverstärkt ($A_{\mathrm{S}}$ $0{,}8$ bis $1{,}0$; $A_{\mathrm{W}}$ $0{,}20$ bis $0{,}30$), Entscheidung offen \\
        \hline
    \end{tabular}
\end{table}

% Herkunft:
% Literaturschlüssel: In 00_literatur/Literaturverzeichnis_BA-Kuhn_Tim_2024-04-28_23-54-47.bib ist
%   Sto18 = Stoffregen, Motorradtechnik (2018) und Sto18 = Stommel et al., FEM zur Berechnung von
%   Kunststoff- und Elastomerbauteilen (2018). Hier wird Sto18 verwendet; litbib.bib (erzeugt aus der
%   BA-Bib) enthält noch keinen Stommel-Eintrag. Schlüssel mit dem Nutzer abgleichen (Schreibregel 3.4
%   nennt [Sto18]). PDF-Seite = Druckseite + 18 (Docstring stommel.py); zitiert sind Druckseiten.
% Absatz 1 (Extraktion): Modul-Docstring 20_code/10_gearcore/scripts/fe_odb_extract.py (Felddatei,
%   Spannungsdatei, Teilesystem, Orientierungssystem W3/W4, Mittelung über Elementknotenwerte);
%   ADR-118 Decision 2 und Amendment 2026-10-11 Punkt 1 (erst Tensor mitteln, dann Invarianten).
% Absatz 2 (Kurven, Normpunkte): src/gearcore/fe/evaluation.py Modul-Docstring, step_summary,
%   StepSummary.complete, curve (abgebrochene Schritte außerhalb), grid_subset (A bis E stets
%   enthalten); scripts/fe_series_analysis.py GLOBAL_QUANTITIES und die Bildunterschrift
%   "Größtwerte über die Zahnbreite; senkrechte Linien an den Normpunkten A bis E", Tabelle Maximum,
%   Ort ab A, A bis E.
% Absatz 3 (Verfahren, Lastgeschichte, Werkstoffstufen): Sto18 S. 313 (Ort der höchsten Beanspruchung),
%   S. 343 bis 344 (isotropes Modell, anisotrop nicht zulässig, elasto-plastisch zulässig), S. 362
%   (6.3.1 linear-elastisch); ADR-118 Decision 3, Amendment 2026-10-13 (knotenweise, wechselnd
%   erkannt), Amendment 2026-10-12 Punkt 6 (W2 neben W1); fatigue/__init__.py und check.py Docstrings.
% Absatz 4 (statischer Nachweis): Sto18 S. 322 (q, f_tau), S. 325 bis 326 (h, Gl. 6.3 und 6.4),
%   S. 344 bis 348 (Nomogramm, A, K), S. 346 und 349 (A = 1,0 und Streckgrenze), S. 349 bis 350
%   (Temperaturübertragung), S. 355 (Gl. 6.15, S = 1), S. 357 (Gl. 6.16, 6.17, f_sigma = 1,3);
%   K als Maximum und Rohkurven von Matscape: ADR-118 Amendments 2026-10-11 (Punkt 3 des ersten,
%   Punkt 1 des dritten Amendments), materials.py Docstring; maßgebender Zustand und a_stat unendlich:
%   check.py Docstring, ADR-118 Prüfung P1-1/P1-2 und Amendment 2026-10-12 Punkt 3.
% Absatz 5 (zyklischer Nachweis): Sto18 S. 360 (Gl. 6.18), S. 362 bis 363 (A_W, Gl. 6.19, fester
%   Wert verstärkt), S. 365 (Haigh, M = sigma_W/sigma_Y), S. 366 bis 368 (Überlastungsfälle, F2),
%   S. 369 bis 371 (Gl. 6.22, k, A_S und A_W bei anderer Temperatur), S. 372 (Gl. 6.23, untere Grenze
%   10^4), S. 328 (weiterer Abfall jenseits 10^7); Begründung F2 als Lastpfad: ADR-118 Amendment
%   2026-10-13; Druckseite, Kappe beidseitig, Lebensdauer: stommel.py und check.py Docstrings,
%   ADR-118 Amendment 2026-10-11 Punkt 4 und Prüfung P1-3/P1-4; Empfindlichkeit F1 bis F4: E2 offen.
% Absatz 6 (komponentenweise, Route principal, a_B,V): Sto18 S. 361 (nicht drehende Hauptrichtungen),
%   S. 368 und 373 (komponentenweise), S. 374 (Gl. 6.24 bis 6.26), S. 332 (Gl. 6.9 komponentenweise);
%   Projektion auf die Richtungen des maßgebenden Zustands: ADR-118 Prüfung P0-2, check.py Docstring;
%   Vergleichswege signed_mises, mmk: check.py Docstring (Seiten der Gl. 6.10 bis 6.12 nicht gelesen,
%   daher Platzhalter).
% Absatz 7 (Klassen): Sto18 S. 346 (Bild 6.18: verstärkt teilkristallin einmalig 0,55 bis 0,70),
%   S. 362 (Bild 6.26: A_W 0,25), S. 323 (Tabelle 6.1 mit Fußnote, q = 0), S. 374 (q ohne
%   Druckkorrektur); duktile Klasse und f_tau 0,577 aus Matscape: ADR-118 Amendment 2026-10-11
%   (Antworten des Nutzers, Punkt 1).
% Absatz 8 (Hill-Erweiterung): fatigue/anisotropy.py Modul-Docstring (Formel, Konstanten aus R_ij,
%   Faserrahmen vorzeichenfrei, Feldvariablen, Invarianten, signed_hill);
%   ADR-118 Amendment E7 Punkte 1, 2, 5, 6b und Prüfung 7 (druckkorrigierte Hauptspannungen
%   zurückgelegt, lineare Interpolation über das Gitter, Hauptrichtungsweg nicht mit Hill);
%   Sto18 S. 345 bis 346 (pauschaler Abminderungsfaktor quer zur Faser, geringe plastische Reserve);
%   Klasse unter Hill: fe_fatigue.py HILL_DEFAULT_CLASS, ADR-118 E7 Punkt 6b; Abaqus-Quelle fehlt in
%   der Bib, daher Platzhalter.
% Absatz 9 (Validierung): ADR-118 Decision 3 und 4 (Funktionen je Gleichung, Buchzahlen als Tests),
%   Amendments 2026-10-11 (Matscape als Oracle der Werkstoffseite: sigma_B,S, sigma_SK, sigma_WK, M;
%   Rohkurven), Amendment 2026-10-14 (E4: S-Life 4.0.0, Einstellungen, 24 Berichte, Knoten 6193 ohne
%   Bericht, Transkription über die Textebene, Hauptspannungen verglichen);
%   80_output/fe/kst_e/frozen_mesh_20layers_bore33_gerechnet/stommel/vergleich_slife.md Kopf und
%   Anmerkungen (Vorgehen, keine Zahlen übernommen); Sto18 S. 375 (Hinweis auf Software [SLife2017]);
%   Zahlenbeispiele S. 346 bis 350.
% Tabelle: ADR-118 (Decision 3, Amendments 2026-10-11, 2026-10-12 Punkt 2 und 3, 2026-10-13),
%   check.py CheckSettings (load_case einmalig, overload_case 2, route principal, required_cycles N_G,
%   safety 1, f_sigma 1,3, unbelasteter Zustand), stommel.py (h = 0 unbelastet), fe_fatigue.py
%   (HILL_DEFAULT_CLASS); Klassenwerte Sto18 S. 346 und 362.
% Bildverweise: fig:nomogramm_qs, fig:haigh_qs_80C aus 80_output/fe/stommel_methode_V2/bildunterschriften.md
%   (Kapitel 2, hier nicht eingebunden).
% Nomenklatur (Symbol, Einheit, Benennung, Quelle); Schlüssel des Größenverzeichnisses in Klammern,
%   "fehlt" = weder in quantities.yaml noch in nomenklatur.yaml gefunden:
% Nomenklatur: sigma_1, MPa, größte Hauptnormalspannung, Sto18 S. 357 (max_principal_stress, gelistet)
% Nomenklatur: sigma_2, MPa, mittlere Hauptnormalspannung, Sto18 S. 325 (fehlt)
% Nomenklatur: sigma_3, MPa, kleinste Hauptnormalspannung, Sto18 S. 357 (min_principal_stress, gelistet)
% Nomenklatur: sigma_V, MPa, Vergleichsspannung nach Gl. (6.16), Sto18 S. 357 (fehlt)
% Nomenklatur: sigma_V,Mises, MPa, von-Mises-Spannung (druckkorrigiert), Sto18 S. 322 und 357 (mises_stress, gelistet)
% Nomenklatur: sigma_V,NH, MPa, Vergleichsspannung der Normalspannungshypothese, Sto18 S. 322 und 357 (fehlt)
% Nomenklatur: sigma_V,Hill, MPa, Hill-Vergleichsspannung im Faserkoordinatensystem, Abaqus 2025 Hill (hill_equivalent_stress, gelistet)
% Nomenklatur: sigma_11 ... sigma_23, MPa, Spannungskomponenten im Faserkoordinatensystem, Abaqus 2025 Hill (fehlt)
% Nomenklatur: R_ij, -, Fließspannungsverhältnisse nach Hill, Abaqus 2025 Hill (hill_stress_ratio, gelistet)
% Nomenklatur: F, G, H, L, M, N, -, Konstanten des Hill-Potentials aus R_ij, Abaqus 2025 Hill (fehlt;
%   M und N kollidieren mit Mittelspannungsempfindlichkeit M und geforderter Lastspielzahl N, Nutzer entscheidet)
% Nomenklatur: q, -, Steuergröße, Sto18 S. 322 (control_variable, im Größenverzeichnis, nicht in nomenklatur.yaml)
% Nomenklatur: f_tau, -, Schubfestigkeitsfaktor, Sto18 S. 322 (shear_strength_factor, im Größenverzeichnis, nicht in nomenklatur.yaml)
% Nomenklatur: f_sigma, -, Druckfestigkeitsfaktor, Sto18 S. 357 (compressive_strength_factor, im Größenverzeichnis, nicht in nomenklatur.yaml)
% Nomenklatur: h, -, Mehraxialitätsgrad nach Hencky, Sto18 S. 325 (multiaxiality_degree, im Größenverzeichnis, nicht in nomenklatur.yaml)
% Nomenklatur: epsilon_V, %, Vergleichsdehnung (zu sigma_V auf der Isotherme abgelesen), Sto18 S. 355 (equivalent_strain, im Größenverzeichnis, nicht in nomenklatur.yaml)
% Nomenklatur: epsilon_zul, %, zulässige Dehnung, Sto18 S. 326 und 355 (permissible_strain, im Größenverzeichnis, nicht in nomenklatur.yaml)
% Nomenklatur: S, -, Sicherheitsfaktor (hier 1), Sto18 S. 355 (safety_factor, im Größenverzeichnis, nicht in nomenklatur.yaml)
% Nomenklatur: A, -, Abminderungsfaktor der Werkstoffklasse und Lastart, Sto18 S. 345 bis 346 (fehlt)
% Nomenklatur: A_S, -, Abminderungsfaktor für einmalige statische Belastung, Sto18 S. 371 (static_reduction_factor, im Größenverzeichnis, nicht in nomenklatur.yaml)
% Nomenklatur: A_W, -, zyklischer Abminderungsfaktor, Sto18 S. 363 (cyclic_reduction_factor, im Größenverzeichnis, nicht in nomenklatur.yaml)
% Nomenklatur: K, MPa, statischer Werkstoffkennwert, Sto18 S. 363 (static_characteristic_value, im Größenverzeichnis, nicht in nomenklatur.yaml)
% Nomenklatur: sigma_stat, MPa, statische Bauteilfestigkeit, Sto18 S. 348 und 372 (fehlt)
% Nomenklatur: sigma_Y, MPa, Streckspannung (statisch), Sto18 S. 368 (fehlt)
% Nomenklatur: a_stat, -, statischer Auslastungsgrad, Sto18 S. 355 (static_utilisation, gelistet)
% Nomenklatur: sigma_B, MPa, Bauteil-Betriebsfestigkeit (Amplitudenwert), Sto18 S. 360 (fehlt; das
%   Größenverzeichnis führt unter static_strength sigma_B als statische Festigkeit nach Gl. (6.7) S. 329, Kollision klären)
% Nomenklatur: sigma_W, MPa, Bauteil-Wechselfestigkeit, Sto18 S. 360 (component_alternating_strength, gelistet)
% Nomenklatur: K_B, -, Betriebsfestigkeitsfaktor, Sto18 S. 360 (fatigue_life_factor, gelistet)
% Nomenklatur: K_M, -, Mittelspannungsfaktor, Sto18 S. 360 (mean_stress_factor, gelistet)
% Nomenklatur: M, -, Mittelspannungsempfindlichkeit, Sto18 S. 368 (mean_stress_sensitivity, gelistet)
% Nomenklatur: sigma_a, MPa, Spannungsamplitude, Sto18 S. 368 (stress_amplitude, gelistet)
% Nomenklatur: sigma_a,i, MPa, Spannungsamplitude der Hauptspannungskomponente i, Sto18 S. 374 (fehlt)
% Nomenklatur: sigma_m, MPa, Mittelspannung, Sto18 S. 368 (mean_stress, gelistet)
% Nomenklatur: R, -, Spannungsverhältnis, Sto18 S. 365 (fehlt)
% Nomenklatur: N, -, geforderte Lastspielzahl, Sto18 S. 371 (required_number_of_cycles, gelistet)
% Nomenklatur: N_G, -, Grenzlastspielzahl 10^7, Sto18 S. 371 (limit_number_of_cycles, gelistet)
% Nomenklatur: k, -, Neigungsexponent der Wöhlerkurve, Sto18 S. 371 (woehler_slope_exponent, gelistet)
% Nomenklatur: a_sigma,i, -, Einzelauslastungsgrad der Hauptspannungskomponente i, Sto18 S. 374 (fehlt)
% Nomenklatur: a_B,V, -, Vergleichsauslastungsgrad (Bauteil-Betriebsfestigkeit), Sto18 S. 374 (cyclic_utilisation, gelistet)
% Nomenklatur: a_B,NH, -, Vergleichsauslastungsgrad nach der Normalspannungshypothese, Sto18 S. 374 (fehlt)
% Nomenklatur: a_B,Mises, -, Vergleichsauslastungsgrad nach von Mises, Sto18 S. 374 (fehlt)
% Nomenklatur: i, Index der Hauptspannungen, Sto18 S. 357 (Index, fehlt in indices)
% Nomenklatur: Abkürzung CONVERSE, Mapping-Programm der Werkstoffkarte (fehlt in abkuerzungen); S-Life Plastics (Programm, keine Abkürzung)
